\chapter{Kết luận}

Khóa luận xuất phát từ khoảng trống căn bản trong các kiến trúc thao tác robot dựa trên mô hình thị giác, ngôn ngữ và hành động đương đại: sự thiếu vắng cơ chế theo dõi trạng thái lịch sử trong các tác vụ dài hạn. Khi thực hiện các thao tác có tính chu kỳ, lặp lại hoặc phụ thuộc danh tính thực thể, các mô hình VLA hiện nay phần lớn vận hành như chính sách Markov thuần túy, chỉ tiếp nhận ảnh quan sát tức thời và câu lệnh ngôn ngữ tĩnh. Khi môi trường trải qua các biến đổi trạng thái không thể phân biệt bằng giác quan hiện tại, chính sách dễ rơi vào bế tắc hoặc lặp lại thao tác vô tận do thiếu cơ sở xác định tiến độ hoàn thành. Trước thách thức đó, khóa luận đề xuất cách tiếp cận mang tính nguyên lý: thay vì lưu trữ toàn bộ chuỗi quan sát thị giác thô gây bùng nổ ngữ cảnh, hệ thống xây dựng kiến trúc bộ nhớ sự kiện rời rạc có bằng chứng vật lý xác thực. Bộ nhớ ghi nhận các biến cố chuyển đổi trạng thái then chốt, biểu diễn quan hệ thực thể và chiếu thành ngữ cảnh ngôn ngữ súc tích dưới ngân sách token cố định. Đóng góp nền tảng của nghiên cứu chứng minh rằng năng lực nhận biết lịch sử và điều phối thao tác dài hạn hoàn toàn có thể tích hợp vào các mô hình nền VLA sẵn có thông qua giao diện ngôn ngữ tự nhiên mà không cần can thiệp cấu trúc mạng nơ-ron cơ sở.

Về mặt lý thuyết và phương pháp luận, khóa luận xác lập ba trụ cột đóng góp chính. Thứ nhất, nghiên cứu hình thức hóa mô hình máy trạng thái hữu hạn cho bộ ghi sự kiện, thiết lập bất biến kiểm tra tính hợp lệ của lịch sử và chứng minh toán học tính chất không nhìn trước của thông tin bộ nhớ, bảo đảm nội dung cung cấp cho chính sách chỉ phản ánh các biến cố đã thực sự diễn ra kèm mốc thời gian khả dụng có bằng chứng xác nhận. Thứ hai, khóa luận đề xuất cơ chế điều phối truy vấn linh hoạt theo sự kiện, giải phóng robot khỏi ràng buộc thực thi cố định đoạn hành động mở đoạn khi môi trường đã chuyển pha thao tác, thiết lập nhịp suy luận thích ứng với động học thực tế. Thứ ba, nghiên cứu tiên phong xây dựng quy trình thực nghiệm đa tầng nghiêm ngặt nhằm tách bạch giữa năng lực sinh hành động cơ học, độ nhạy đối với ngữ cảnh ngôn ngữ và khả năng bảo toàn tri thức sau tinh chỉnh. Quy trình này mở ra góc nhìn mới về sự phân tách trách nhiệm trong thiết kế robot thông minh: quyết định kết thúc chu kỳ thao tác nên được nâng đỡ bởi cơ chế ngoại vi dựa trên bằng chứng vật lý vững chắc, thay vì đặt trọn gánh nặng suy luận lên mô hình nền vốn nhạy cảm với sự trôi dạt phân phối.

Bên cạnh các kết quả đạt được, khóa luận nhìn nhận thẳng thắn những hạn chế nội tại cần khắc phục. Trước hết, toàn bộ các kiểm chứng mới chỉ tiến hành trong môi trường mô phỏng vật lý với số lượng tác vụ và trạng thái khởi tạo có giới hạn, chưa phản ánh hết các nhiễu loạn động lực học, sai số cơ khí và biến đổi ánh sáng trong thế giới thực. Thêm vào đó, việc phát hiện biến cố chuyển tiếp vẫn dựa vào sự kết hợp giữa mô hình thị giác mở và cảm biến trạng thái nội tại với những ngưỡng quy tắc tĩnh, khiến hệ thống dễ bị ảnh hưởng bởi độ trễ tính toán của mạng thị giác lớn cũng như tình trạng che khuất góc nhìn camera. Về bản chất mô hình hóa, các thử nghiệm can thiệp chỉ ra rằng chính sách VLA hiện tại tuy nhạy bén với sự hiện diện của đoạn văn bản chỉ dẫn nhưng dường như chưa thực sự thấu hiểu ngữ nghĩa sâu của dữ kiện lịch sử để tự động điều chỉnh chuỗi hành động tinh tế. Sự phân hóa về mức độ suy giảm kỹ năng cơ sở giữa các kiến trúc mô hình nền khác nhau cũng cho thấy tinh chỉnh tham số phục vụ bộ nhớ vẫn tiềm ẩn rủi ro quên thảm khốc nếu thiếu kỹ thuật điều hòa phân phối dữ liệu chặt chẽ hơn.

Từ những hạn chế trên, nghiên cứu mở ra các hướng phát triển khoa học giàu triển vọng cho cộng đồng robot học và mô hình ngôn ngữ, thị giác, hành động. Hướng đi ưu tiên hàng đầu là chuyển giao và kiểm chứng toàn diện kiến trúc bộ nhớ sự kiện lên các cánh tay robot vật lý thực tế, nơi dữ liệu xúc giác, lực tiếp xúc và hình ảnh nổi cung cấp nguồn bằng chứng đa phương thức phong phú với độ trễ thấp hơn luồng thị giác đơn thuần. Về mặt kiến trúc, thay vì phụ thuộc ngưỡng quy tắc tiền định, nghiên cứu tiếp theo có thể hướng tới tối ưu hóa đầu dự đoán trạng thái dừng nội tại hoặc cơ chế chú ý theo sự kiện huấn luyện liên kết trực tiếp với biểu diễn tiềm ẩn của mô hình nền, cho phép chính sách tự chủ hoàn toàn trong việc ngắt chu kỳ hành động. Bên cạnh đó, việc xây dựng cơ chế suy luận bất đồng bộ giữa luồng điều khiển vận động tần số cao và luồng cập nhật bộ nhớ nhận thức tần số thấp sẽ giúp triệt tiêu độ trễ điều khiển, bảo đảm an toàn chuyển động cho robot trong môi trường tương tác biến đổi liên tục. Cuối cùng, mở rộng khung biểu diễn sang kịch bản thao tác hợp tác đa robot và các bài toán suy luận nhân quả trong không gian mở sẽ là bước tiến quyết định đưa hệ thống VLA tiệm cận năng lực hành động tự chủ.
